%!TEX program = lualatex
%!TEX encoding = UTF-8 Unicode  
%!TEX root = chapter15.tex  
\documentclass[main.tex]{subfiles} 
\begin{document} 
%----------------------------------------------------------------------------------------
%	CHAPTER 15
%---------------------------------------------------------------------------------------- 
\cleardoublepage 
\loadgeometry{marginlayout} 
\pagestyle{scrheadings} % Print headers again  
\KOMAoptions{
headwidth=textwithmarginpar,
headsepline=true,
headsepline= 0.5pt,
footwidth=textwithmarginpar,
}   
\onehalfspacing 
\setcounter{chapter}{14} 


\chapter{Factoriser pour résoudre des équations produit et quotient}


%%----------------------------------------------------------------------------------------
%%	   
%%----------------------------------------------------------------------------------------
% 
%
%\newpage 
%\loadgeometry{marginlayout}  
%\pagestyle{scrheadings} % Print headers again  
%\KOMAoptions{
%	headwidth=textwithmarginpar,
%	headsepline=true,
%	headsepline= 0.5pt,
%	footwidth=textwithmarginpar,
%} 
%\onehalfspacing
\section{Équations produit nul} 
\begin{theorem}[produit nul] Si $AB=0$ alors $A=0$ ou $B=0$
\end{theorem}
\begin{proof}On suppose $AB=0$. Par disjonction de cas :
\begin{enumerate}[leftmargin=*,label=\textbullet]
	\item Soit $A=0$. Rien à démontrer.
	\item Soit $ A\neq 0$. On va démontrer que $B=0$.
	
	En effet :  $A$ admet un inverse $\frac{1}{A}$ et son inverse est non nul.\\
	\noindent \parbox[position=t]{0.85\linewidth}{	\color{white}
		$\begin{WithArrows}
			AB &=0\Arrow{$\times \frac{1}{A}$}\\
		 \frac{1}{A}\times A \times B&=\rule{0pt}{3em} \frac{1}{A}\times  0\Arrow{$\frac{1}{A}\times A=1$}\\
			B& =0 \rule{0pt}{2em}  
		\end{WithArrows}$ 	
	} 
\end{enumerate} 
\end{proof}
\begin{eBox}
	La forme développée simplifiée réduite d'une expression du second degré  $f(x)=ax^2+bx+c$ n'est pas adaptée dans la recherche des zéros de $f(x)=0$. 
\end{eBox}
\begin{eBox} 
	Trouver $r_1$ et $r_2$ tel que pour tout : $x\in\mathbb{R}$ on a $f(x)=a(x-r_1)(x-r_2)$ (forme factorisée), c'est trouver les zéros de $f$ : $f(r_1)=0$ et $f(r_2)=0$. 
\end{eBox}
\begin{eBox} 
	Si une expression du second degré ne s'annulle pas, alors elle n'est pas factorisable sous la forme $a(x-r_1)(x-r_2)$ avec $r_1$ et $r_2\in\mathbb{R}$ !
\end{eBox}
\begin{example}
Les expressions $A(x)=5x^2+3$ et $B(x)=-2(3x+5)^2-1$ ne sont pas factorisables dans $\mathbb{R}$. En effet  pour tout $x\in\mathbb{R}$ on a $A(x) \geqslant 3$ et $B(x) \leqslant -1$.  
%$C(x)=5x^2-3$ est factorisable.% Les solutions de $C(x)=0$ sont $\sqrt{\tfrac{3}{5}}$ et $-\sqrt{\tfrac{3}{5}}$ et $C(x)=5(x-\sqrt{\tfrac{3}{5}})(x+\sqrt{\tfrac{3}{5}})$
\end{example}
%----------------------------------------------------------------------------------------
%	EXERCICES
%----------------------------------------------------------------------------------------

\newpage 
\loadgeometry{widelayout}  
\pagestyle{scrheadings} % Print headers again  
\KOMAoptions{
headwidth=textwithmarginpar,
headsepline=true,
headsepline= 0.5pt,
footwidth=textwithmarginpar,
} 
\onehalfspacing
\subsection{Exercices : équations produit nul, forme factorisée pour résoudre}



\begin{example}  Factoriser/réduire le membre de gauche et résoudre dans $\mathbb{R}$ les équations inconnue $x$: 

\noindent  \parbox{0.24\linewidth}{	
$\begin{WithArrows}
4x^2-9x&=0 \\
&  \rule{0pt}{1.3em} \\
&   \rule{0pt}{1.3em}  \\
&    \rule{0pt}{1.3em}  \\
&   \rule{0pt}{1.0em}  
\end{WithArrows}$ 
}  \hfill \parbox{0.24\linewidth}{	
$\begin{WithArrows}
4x^2-9&=0 \\
&  \rule{0pt}{1.3em} \\
&   \rule{0pt}{1.3em}  \\
&   \rule{0pt}{1.3em}  \\
&   \rule{0pt}{1.0em}  
\end{WithArrows}$ 
}  \hfill \parbox{0.24\linewidth}{	
$\begin{WithArrows}
	4x^2+9&=0 \\
	&  \rule{0pt}{1.3em} \\
	&   \rule{0pt}{1.3em}  \\
	&   \rule{0pt}{1.3em}  \\
	&   \rule{0pt}{1.0em}  
\end{WithArrows}$ 
}    \hfill \parbox{0.24\linewidth}{	
$\begin{WithArrows}
(-2x-7)^{15}&=0 \\
&  \rule{0pt}{1.3em} \\
&   \rule{0pt}{1.3em}  \\
&   \rule{0pt}{1.3em}  \\
&   \rule{0pt}{1.0em}   
\end{WithArrows}$ 
}      



\noindent \parbox{0.31\linewidth}{	
$\begin{WithArrows}
x^2&=12x-36 \\
&  \rule{0pt}{1.3em} \\
&  \rule{0pt}{1.3em}  \\
&   \rule{0pt}{1.3em}  \\
&   \rule{0pt}{1.3em}    \\
&   \rule{0pt}{1.0em}  
\end{WithArrows}$ 	
}\hfill \parbox{0.31\linewidth}{	
$\begin{WithArrows}
9x^2+24x+4&=0 \\
&  \rule{0pt}{1.3em} \\
&   \rule{0pt}{1.3em}  \\
&    \rule{0pt}{1.3em}  \\
&   \rule{0pt}{1.3em}   \\
&   \rule{0pt}{1.0em}  
\end{WithArrows}$ 
}  \hfill \parbox{0.31\linewidth}{	
$\begin{WithArrows}
(x+1)(2x+3)+5(x+1)&=0 \\
&  \rule{0pt}{1.3em} \\
&   \rule{0pt}{1.3em}  \\
&   \rule{0pt}{1.3em}  \\
&  \rule{0pt}{1.3em}   \\
&   \rule{0pt}{1.0em}  
\end{WithArrows}$ 
}    
\end{example} 
\vspace{-0.5em}
\begin{exercise}\label{chapter07:exercice:equationproduitA} Résoudre dans $\mathbb{R}$ les équations suivantes, inconnue $x$. 
\vspace{-0.5em}
\begin{multicols}{3}
\begin{enumerate}[leftmargin=*,label={$(E_{\arabic*})$}]
\item $(5x-10)(8x+5)=0$
\item $(8x-1)(7x-4)=0$
\item $(8x-1)+(7x-4)=0$ 
\item $5x^2+2x=0$
\item $ 9x^2-4=0$
\item $36x^2+100=0$
\item $x^2+10x+25=0$ 
\item $-3x^2=5x$
\item $x^2=2x-1$
\item $x^2+4x=-4$
\item $(x-3)^2(2x-1)^4=0$  
\item $x^2-6x=-9$
\end{enumerate}
\end{multicols} 
\end{exercise}  
\vspace{-0.5em}
\begin{exercise}\label{chapter07:exercice:equationproduitB} Mêmes consignes \vspace{-0.5em}
\begin{multicols}{2}
\begin{enumerate}[leftmargin=*,label={$(E_{\arabic*})$}]
\item $(3x+2)(8x+5)+(3x+2)(2x+3)=0$ 
\item $(3x-1)^2+(3x-1)(5x-4)=0$ 
\item $(3x-1)^2-(5x-4)^2=0$ 
%			\item  $0=(3x+2)(8x+5)+(3x+2)(2x+3)$
\item $(2x-5)(5x-4)=(2x-5)(8x-1)$ 
\item $(6x-4)(-3x+2)=(10x-2)(6x-4)$  
\item $(2x+3)^2=(5x+4)^2$ 
\end{enumerate}
\end{multicols} 
\end{exercise}
\vspace{-0.5em} 
\begin{exercise}[une cubique] Soit $f$ et $g$ définies sur $\mathbb{R}$ par $f(x)=x^3-7x+6$ et $g(x)=x^2+2x-3$.
\begin{enumerate}[label=\arabic*), leftmargin=*]
	\item Vérifiez que $g(-3)=0$.
	\item Déterminer $a$ et $b$ tel que  pour tout $x\in\mathbb{R}$ on a $g(x)=(x+3)(ax+b)$.
	\item Vérifiez que $f(2)=0$.
	\item Déterminer $A$, $B$ et $C$ tel que pour tout $x\in\mathbb{R}$ on a $f(x)=(x-2)(Ax^2+Bx+C)$
	\item En déduire toutes les solutions de l'équation $f(x)=0$.
\end{enumerate}
\end{exercise}

\begin{exercise}[Un grand classique : choisir la forme algébrique la plus adaptée]

On considère l'expression définie pour tout $x$ appartenant à $\mathbb{R}$ par $A(x) =(x+3)^2+2(x+1)(x+3)$.
\begin{enumerate}[label=\arabic*), leftmargin=*]
\item Développer, réduire et ordonner $A(x)$  (forme développée). 
\item Factoriser $A(x)$ et montrer que $A(x)=(x+3)(3x+5)$ (forme factorisée).
\item Calculer $A(0)$.
\item En utilisant la forme développée donner la valeur exacte de $A(\sqrt{2})$. 
\item En utilisant la forme factorisée, résoudre l'équation $A(x)=0$.
\item Utiliser la forme la plus adaptée pour déterminer la valeur des deux solutions de $A(x)=15$.
\end{enumerate}
\end{exercise}

\begin{exercise}[rebelotte] Soit la fonction $B$ définie sur $\mathbb{R}$ par    $B(x) =(2x-4)^2-3(x-2)(x+5)$.
\begin{enumerate}[label=\arabic*), leftmargin=*]
\item Développer, réduire et ordonner $B(x)$ (forme développée). 
\item Montrer que pour tout $x\in\mathbb{R}$,  $B(x)=(x-23)(x-2)$ (forme factorisée).
\item En utilisant la forme la plus adaptée calculer $B(0)$ et $B(-\sqrt{2})$.  
\item En utilisant la forme la plus adaptée résoudre dans $\mathbb{R}$ l'équation $B(x)=0$ d'inconnue $x$.  
\item En utilisant la forme la plus adaptée résoudre dans $\mathbb{R}$ l'équation $B(x)=46$ d'inconnue $x$.  
\item Montrer que pour tout $x\in\mathbb{R}$, $B(x)=(x-\numprint{12.5})^2- \numprint{110.25}$ (forme canonique).
\item En utilisant la forme canonique résoudre dans $\mathbb{R}$ l'équation $B(x)=- \numprint{110.25}$ d'inconnue $x$. 
\item En utilisant la forme canonique résoudre dans $\mathbb{R}$ l'équation $B(x)=\numprint{10.75}$ d'inconnue $x$. 
\end{enumerate}
\end{exercise}

\begin{exercise}[un dernier]Soit la fonction $C$ définie sur $\mathbb{R}$ par $C(x) =(3x+5)^2-(2x+3)^2$.
\begin{enumerate}[label=\arabic*), leftmargin=*]
\item Développer, réduire et ordonner $C(x)$. 
\item Factoriser $C(x)$ à l'aide d'une identité remarquable et montrer que $C(x)=(5x+8)(x+2)$.
\item Calculer $C(0)$.
%\item En utilisant la forme développée donner la valeur exacte de $C(-\sqrt{2})$. 
\item En utilisant la forme factorisée, résoudre l'équation $C(x)=0$.
\item L'équation  $C(x)=16$ a deux solutions. En utilisant la forme la plus adaptée, déterminer les.
\end{enumerate}
\end{exercise}

\begin{exercise} \hfill 
	
	
	\noindent\parbox{0.30\linewidth}{\centering 
		\begin{tikzpicture}[scale=0.7]
			\tkzDefPoint(-2,-1.5){A}
			\tkzDefPoint(2.5,-1.5){B}
			\tkzDefPoint(2.5,1){C}
			\tkzDefPoint(-2,1){D}
			\tkzDrawSegments[line width=1.2pt](A,B B,C C,D D,A)
			\tkzLabelSegment[below](A,B){\small$x+4$}
			\tkzLabelSegment[right](B,C){\small$x$}
			\tkzLabelSegment[below](A,C){\small Aire = \SI{60}{\centi\meter^2}}
			%		\tkzLabelSegment[above](D,C){\scriptsize$4x+2$}
			\tkzMarkRightAngle[  size=0.25](B,A,D) 
			\tkzMarkRightAngle[  size=0.25](C,B,A) 
			\tkzMarkRightAngle[  size=0.25](D,C,B) 
			\tkzMarkRightAngle[  size=0.25](A,D,C) 
		\end{tikzpicture} 
	}\hfill \parbox{0.65\linewidth}{ 
		\begin{enumerate}[label=\arabic*), leftmargin=*]
			\item  Montrer que $x$ vérifie $x^2+4x-60=0$. 
			\item Factoriser par essai-erreur et identifier les solutions admissibles.
		\end{enumerate} 
	}   
\end{exercise}
\begin{exercise} \hfill 
	
	\noindent\parbox{0.30\linewidth}{\centering 
		\begin{tikzpicture}[scale=0.7]
			\tkzDefPoint(-2,-1.5){A}
			\tkzDefPoint(2.5,-1.5){B}
			\tkzDefPoint(2.5,1){C}
			\tkzDefPoint(-2,1){D}
			\tkzDrawSegments[line width=1.2pt](A,B B,C C,D D,A)
			\tkzLabelSegment[below](A,B){\small$x+5$}
			\tkzLabelSegment[below, sloped](B,C){\small$x-4$}
			\tkzLabelSegment[below](A,C){\small Aire = \SI{36}{\centi\meter^2}}
			%		\tkzLabelSegment[above](D,C){\scriptsize$4x+2$}
			\tkzMarkRightAngle[  size=0.25](B,A,D) 
			\tkzMarkRightAngle[  size=0.25](C,B,A) 
			\tkzMarkRightAngle[  size=0.25](D,C,B) 
			\tkzMarkRightAngle[  size=0.25](A,D,C) 
		\end{tikzpicture}  
	}\hfill \parbox{0.65\linewidth}{ 
		\begin{enumerate}[label=\arabic*), leftmargin=*] 
			\item  Montrer que $x$ vérifie $x^2+x-56=0$. 
			\item Factoriser par essai-erreur et identifier les solutions admissibles.
		\end{enumerate} 
	}   
\end{exercise}
\begin{exercise} \hfill 


\noindent\parbox{0.30\linewidth}{\centering 
	\begin{tikzpicture}[scale=.75]
		\tkzSetUpPoint[size=2, fill=black]
		\tkzfctset{tan style/.style={->,>=latex, color=red, line width=1.2pt}}   
		\tkzSetUpLine[line width= 1.2pt, add=.5 and .5]  
		%			\tkzInit[xmin=-1.0,ymin=-1.0,xmax=6.5,ymax=2.8]
		%			\tkzClip
		\tkzDefPoint(0,0){A} 
		\tkzDefPoint(4.5,0){B}
		\tkzDefPoint(4.5,2){C}
		\tkzDefPoint(0,2){D} 
		\tkzDrawPolygon[line width=1.2pt](A,B,C,D)  
		
		\tkzDrawSegment[dim={\small$x+2$,-15pt,}](A,B)
		\tkzLabelSegment[above,font=\small](A,C){$58$}
		\tkzDrawSegments[dashed, line width=1.2pt](A,C)  
		\tkzDrawSegment[dim={\small$x$,15pt,},dashed](A,D)  
%		\tkzMarkRightAngle[  size=0.25](B,A,F) 
	\end{tikzpicture}  
}\hfill \parbox{0.65\linewidth}{ 
	La figure ci-contre est un   rectangle. 
	\begin{enumerate}[label=\arabic*), leftmargin=*]
		\item Montrer que $x$ vérifie  $x^2+2x-1680=0$.
		\item Développer $(x-40)(x+42)$.
		\item En déduire le(s) valeur(s) possibles de $x$.
	\end{enumerate}
}  

\end{exercise}	
\begin{exercise} \hfill 
	
	\noindent\parbox{0.30\linewidth}{\centering 
		\begin{tikzpicture}[scale=.8]
			\tkzSetUpPoint[size=2, fill=black]
			\tkzfctset{tan style/.style={->,>=latex, color=red, line width=1.2pt}}   
			\tkzSetUpLine[line width= 1.2pt, add=.5 and .5]  
			%			\tkzInit[xmin=-1.0,ymin=-1.0,xmax=6.5,ymax=2.8]
			%			\tkzClip
			\tkzDefPoint(0,0){A} 
			\tkzDefPoint(5,0){B}
			\tkzDefPoint(0.8,2.25){D}
			\tkzDefPoint(3.25,2.25){C}
			\tkzDrawPolygon[line width=1.2pt](A,B,C,D)
			\tkzDefPoint(2,0){I}
			\tkzDefPoint(2,2.25){J}
			\tkzDrawSegments[dashed,line width=0.4pt](I,J A,C)
			
			\tkzDrawSegment[dim={\small$3x-2$,-15pt,}](A,B))
			\tkzDrawSegment[dim={\small$x+5$,-15pt,}](C,D))
			\tkzDrawSegment[dim={\small$2x-3$,-75pt,},dashed](I,J))
			%			\tkzDefTriangle[two angles = 27 and 90](A,C)\tkzGetPoint{B}
			%			\tkzDrawPolygon[line width=1.2pt](A,B,C)
			%			% 			\tkzLabelSegment[below](A,C){$2\sqrt{xy}$}
			%			\tkzLabelSegment[below](A,C){$x$}
			%			\tkzLabelSegment[right](B,C){$x-4$}
			%			\tkzLabelSegment[above left](A,B){$x+4$}
			%			%    	\tkzDrawPoints[color=red](A,B,C)
			%			%   	 \tkzLabelPoints[below](A,C)
			%			%    		\tkzLabelPoints[above right](B)
			\tkzMarkRightAngle[  size=0.25](B,I,J)
			\tkzMarkRightAngle[  size=0.25](I,J,C)
		\end{tikzpicture}  
	}\hfill \parbox{0.65\linewidth}{  
		L'aire du trapèze ci-dessous est de \SI{133}{\centi\meter^2}. Les longueurs sont en \si{\centi\meter}.
		\begin{enumerate}[label=\arabic*), leftmargin=*]
			\item Montrer que $x$ vérifie  $8x^2-6x-275=0$.
			\item Développer $(4x-25)(2x+11)$.
			\item Résoudre en $x$ et préciser la solution admissible.
		\end{enumerate}
	}  
	
\end{exercise}	
\begin{exercise} \hfill 

\noindent\parbox{0.30\linewidth}{\centering 
\begin{tikzpicture}[scale=.75]
\tkzSetUpPoint[size=2, fill=black]
\tkzfctset{tan style/.style={->,>=latex, color=red, line width=1.2pt}}   
\tkzSetUpLine[line width= 1.2pt, add=.5 and .5]  
%			\tkzInit[xmin=-1.0,ymin=-1.0,xmax=6.5,ymax=2.8]
%			\tkzClip
\tkzDefPoint(0,0){A} 
\tkzDefPoint(6,0){B}
\tkzDefPoint(6,2.05){C}
\tkzDefPoint(2,2.05){D}
\tkzDefPoint(2,3.2){E}
\tkzDefPoint(0,3.2){F}
\tkzDrawPolygon[line width=1.2pt](A,B,C,D,E,F)  

\tkzDrawSegment[dim={\small$10$,-10pt,}](A,B))
\tkzDrawSegment[dim={\small$3x+1$,-20pt,}](B,C))
\tkzDrawSegment[dim={\small$9$,15pt,},dashed](A,F))
\tkzDrawSegment[dim={\small$x$,-10pt,},dashed](E,F))
%			\tkzDefTriangle[two angles = 27 and 90](A,C)\tkzGetPoint{B}
%			\tkzDrawPolygon[line width=1.2pt](A,B,C)
%			% 			\tkzLabelSegment[below](A,C){$2\sqrt{xy}$}
%			\tkzLabelSegment[below](A,C){$x$}
%			\tkzLabelSegment[right](B,C){$x-4$}
%			\tkzLabelSegment[above left](A,B){$x+4$}
%			%    	\tkzDrawPoints[color=red](A,B,C)
%			%   	 \tkzLabelPoints[below](A,C)
%			%    		\tkzLabelPoints[above right](B) 
\tkzMarkRightAngle[  size=0.25](B,A,F)
\tkzMarkRightAngle[  size=0.25](C,B,A)
\tkzMarkRightAngle[  size=0.25](D,C,B)
\tkzMarkRightAngle[  size=0.25](C,D,E)
\tkzMarkRightAngle[  size=0.25](F,E,D)
\tkzMarkRightAngle[  size=0.25](A,F,E)
\end{tikzpicture}  
}\hfill \parbox{0.65\linewidth}{ 
L'aire de la figure en L ci-contre est   \SI{65}{\centi\meter^2}. Les longueurs  sont en \si{\centi\meter}. 
\begin{enumerate}[label=\arabic*), leftmargin=*]
\item Montrer que $x$ vérifie    $3x^2-38x+55=0$.
\item Développer $(3x-5)(x-11)$.
\item Résoudre en $x$ et préciser quelles solutions sont admissibles.
\end{enumerate}
}    
\end{exercise}


%----------------------------------------------------------------------------------------
%	   
%----------------------------------------------------------------------------------------


\newpage 
\loadgeometry{marginlayout}  
\pagestyle{scrheadings} % Print headers again  
\KOMAoptions{
	headwidth=textwithmarginpar,
	headsepline=true,
	headsepline= 0.5pt,
	footwidth=textwithmarginpar,
} 
\onehalfspacing
 \section{Fractions algébriques}
 \begin{eBox}
 	Pour certaines expressions de $x$, il convient avant de procéder à leur manipulation de préciser leur \textbf{domaine de définition}. Il s'agit des valeurs de la variable $x$ pour lesquelles l'expression a un sens, en particulier les \textbf{dénominateurs} ne doivent pas s'annuler.   
 \end{eBox}
 \begin{example}Préciser le domaine de définition des expressions suivantes :
 	
 	\noindent  \parbox{0.45\linewidth}{	
 		$\begin{WithArrows}
 			A(x)&=\frac{3x-15}{6x-18} \\
 			\textbf{Valeur Interdite : }	&   6x-18=0 \rule{0pt}{1.3em} \\
 			&  6x = 18 \rule{0pt}{1.3em}  \\
 			&  x=3 \rule{0pt}{1.3em}  
 		\end{WithArrows}$ \\
 		Le domaine de $A$ est $D=\mathbb{R}\setminus \{3\}$.
 	}  \hfill \parbox{0.45\linewidth}{	
 		$\begin{WithArrows}
 			B(x)&=\frac{x-3}{x^2-2x} \\ 
 			\textbf{V.I. : }	& x^2-2x=0  \rule{0pt}{1.3em}  \\
 			& \color{white} x(x-2) = 0 \rule{0pt}{1.3em}  \\
 			& \color{white} x=0 \; \text{ou} \;x= 2 \rule{0pt}{1.3em}    
 		\end{WithArrows}$ \\
 		Le domaine de $B$ est \color{white} $D=\mathbb{R}\setminus \{0~;~2\}$.
 	}     
 \end{example} 
 
 \section{Équations rationnelles}
 
 
 \begin{theorem}[Équation quotient nul] Si $\frac{A}{B}=0$ alors $A=0$ et $B\neq0$
 \end{theorem}
 \begin{proof}\color{white} Comme $\frac{A}{B}=A\times \frac{1}{B}$, $B$ doit avoir un inverse. D'où $B\neq0$. 
 	
 	\noindent \parbox{0.85\linewidth}{	
 		$\begin{WithArrows}
 			\frac{A}{B} &=0 \\
 			A\times \frac{1}{B}&=  0\Arrow{$\times B$}\\
 			A\times \frac{1}{B}\times B & = 0\times B\\
 			A & = 0
 		\end{WithArrows}$ 	
 	} 
 \end{proof}
 \begin{corollaryT} Si $\frac{A}{B}=\frac{C}{D}$ alors $AD=BC$, $C\neq 0$ et $D\neq0$. 
	\end{corollaryT}
%----------------------------------------------------------------------------------------
%	EXERCICES
%----------------------------------------------------------------------------------------

\newpage 
\loadgeometry{widelayout}  
\pagestyle{scrheadings} % Print headers again  
\KOMAoptions{
headwidth=textwithmarginpar,
headsepline=true,
headsepline= 0.5pt,
footwidth=textwithmarginpar,
} 
\onehalfspacing

\subsection{Exercices : fractions algébriques}
\begin{eBox}
\textbf{Simplification de fractions}	Pour $a$ et $c\neq 0$ on a  $\frac{ab}{ac}=\frac{b}{c}$
\end{eBox}
\begin{example}[Simplifier]  les expressions suivantes en trouvant un facteur commun :

\noindent\hfill\parbox{0.43\linewidth}{	
$\begin{WithArrows}
A(x)&= \frac{5x^2+2x}{3x} \rule{0pt}{1.8em} \\ % 
&= \frac{x(5x+2)}{3x}\rule{0pt}{1.8em}  \\  
& =\frac{5x+2}{3} \rule{0pt}{1.8em}  
\end{WithArrows}$ 	
$\begin{WithArrows}
\textbf{V.I. : }& 3x =0 \\
 & D=\mathbb{R}\setminus\{3\}  
\end{WithArrows}$ 	
}    \hfill\parbox{0.55\linewidth}{	
$\begin{WithArrows}
B(x)&= \frac{x^2+4x+4}{x^2-4}  \rule{0pt}{1.8em} \\ 
&= \frac{(x+2)^2}{(x-2)(x+2)}\rule{0pt}{1.8em}  \\  
& =\frac{x+2}{x-2}\rule{0pt}{1.8em}  
\end{WithArrows}$
$\begin{WithArrows}
\textbf{V.I. : }& x^2-4 =0 \\
& (x-2)(x+2) =0 \\
& x =2 \text{ ou } x=-2 \\
 & D=\mathbb{R}\setminus\{2~;~-2\}  
\end{WithArrows}$
}   
\end{example} 
%\parbox{0.31\linewidth}{	
%	$\begin{WithArrows}
%		A(x)& =\frac{14x+8}{2x^2+2} \rule{0pt}{1.8em} \\
%		\textbf{V.I. : }& 2x^2+2 =0 \\
%		\textbf{Domaine est }& D=\mathbb{R} \\
%		A(x)&= \frac{2(7x+4)}{2(x^2+1)}\rule{0pt}{1.8em} \\
%		& =\frac{7x+4}{x^2+1}\rule{0pt}{1.8em}     
%	\end{WithArrows}$ 	
%} \hfill
%\vfill 
% \hfill\parbox{0.24\linewidth}{	
%	$\begin{WithArrows}
%		C&= \frac{8+4x}{6x+12}  \rule{0pt}{1.8em} \\ 
%		&=\rule{0pt}{1.8em}  \\  
%		& =\rule{0pt}{1.8em}   \\
%		& =\rule{0pt}{1.8em}   
%	\end{WithArrows}$
\begin{exercise} \label{chapitre02:fractionalgebriques:simplification} Pour chaque fraction, déterminer le domaine $D$ et simplifier la pour $x\in D$. :\hfill \vspace{-0.5em}
\begin{multicols}{4} 
\begin{enumerate}[label={$\Alph*(x)=$\rule{0pt}{1.8em}}, leftmargin=*] 
%\item $\frac{2x-8}{12x-4}$
\item $\frac{2x^2+8x}{4x}$ 
\item $\frac{5x-15}{6x-18}$
\item $\frac{x^2-9}{x+3}$
\item $\frac{x-2}{x^2-2x}$  
\item $\frac{x^2-10x+25}{x-5}$
%\item $\frac{x^2+4x^2}{x^2-x}$
%\item $\frac{x^3+4x^2}{6x^2+4x}$
\item $\frac{x+5}{x^2-25}$ 
\item $\frac{x^2+x}{x^2-1}$  	
%\item $\frac{x^2+x^2}{3x^3}$
\item $\frac{9x^2-12x+4}{9x^2-4}$  
\end{enumerate}
\end{multicols}
\end{exercise} 
\begin{eBox}
	\textbf{Sommes de fractions}	Pour $b$ et $d\neq 0$ on a  $\frac{a}{b}+\frac{c}{d}=\frac{ad}{bd}+\frac{bc}{bd}=\frac{ad+bc}{bd}$
\end{eBox}
\begin{example} Préciser le domaine et ramener au même dénominateur les expressions suivantes.

\noindent\hfill\parbox{0.45\linewidth}{	
$\begin{WithArrows}
A& = \frac{2}{x+1}-1 \rule{0pt}{1.6em} \\
& = \frac{2}{x+1} - \frac{x+1}{x+1}\rule{0pt}{1.6em} \\
& = \frac{2-(x+1)}{x+1} \rule{0pt}{1.6em}   \\  
& = \frac{1-x}{x+1}  \rule{0pt}{1.6em}  
\end{WithArrows}$ 	 	
$\begin{WithArrows}
\textbf{V.I. : }& x+1 =0 \\
& D=\mathbb{R}\setminus\{-1\}  
\end{WithArrows}$ 	
}  \hfill\parbox{0.52\linewidth}{	
$\begin{WithArrows}
B&=  \frac{1}{x}+\frac{2}{x-1}  \rule{0pt}{1.6em} \\ 
& =\frac{1(x-1)}{x(x-1)}+\frac{2x}{(x-1)x} \rule{0pt}{1.6em}  \\  
& =\frac{x-1 + 2x}{(x-1)x} \rule{0pt}{1.6em}  \\  
& =\frac{3x-1}{(x-1)x}  \rule{0pt}{1.6em}   
\end{WithArrows}$ 
$\begin{WithArrows}
	\textbf{V.I. : }& x =0 \text{ et } x-1=0\\
	& D=\mathbb{R}\setminus\{0~;~1\}  
\end{WithArrows}$ 	
}    
\end{example}   
\begin{exercise} \label{chapitre02:fractionalgebriques:memedenominateur}  
Même consignes :\vspace{-0.75em}
\begin{multicols}{3}
\begin{enumerate}[label={$\Alph*(x)=$\rule{0pt}{1.8em}}, leftmargin=*] 
\item $\frac{2}{x+3}+\frac{2}{x-3}$
\item $\frac{4}{x}+\frac{x-1}{3x-5}$
\item $\frac{3}{5x}+\frac{11x}{x+1}$
\item $\frac{3x-2}{5x-3}-\frac{2-3x}{7x-2}$  
\item $\frac{5}{2x-1}-1$ 
\item $2x+1-\frac{2}{2x+1}$
\item $\frac{2}{x-2}-\frac{5}{x+3}$
\item $\frac{1}{x}-\frac{3}{x-5}$
\item $\frac{1}{2x-7} -\frac{1}{2x+9}$
%\item $\frac{3x+4}{2x-9}\times \frac{4x^2-36x+81}{9x^2-16}$
%\item $\frac{x^2-2x+1}{(23x-5)^2}\times \frac{46x-10}{1-x}$
%%\item $\frac{-2}{x+3}+\frac{2}{x-5}$
%%\item $\frac{1}{2(x-1)}+\frac{1}{6(x+1)}$
%%\item $\frac{-2}{3(x+3)}+\frac{1}{x+2}$
%%\item $\frac{3}{4(x-3)}-\frac{5}{8(x-4)}$ 
\end{enumerate}
\end{multicols}  
\end{exercise}
 \begin{eBox}
 	\textbf{Produit et quotient de fractions}	Pour $a, b, c,$ et $d\neq 0$ on a  $\frac{a}{b}\times\frac{c}{d}=\frac{ac}{bd}$ \hfill et \hfill $\frac{a}{b}\divisionsymbol\frac{c}{d}=\frac{a}{b}\times\frac{d}{c}$
 \end{eBox}
\begin{exercise} \label{chapitre02:fractionalgebriques:produit}  
	Même consignes :\vspace{-0.5em}
	\begin{multicols}{3}
		\begin{enumerate}[label={$\Alph*(x)=$\rule{0pt}{1.8em}}, leftmargin=*]  
			\item $\frac{3x+4}{2x-9}\times \frac{4x^2-36x+81}{9x^2-16}$
			\item $\frac{x^2-2x+1}{(23x-5)^2}\times \frac{46x-10}{1-x}$
			\item $\frac{1}{\frac{x+2}{x-5}}  (x+5)$ 
		\end{enumerate}
	\end{multicols}  
\end{exercise}



%----------------------------------------------------------------------------------------
%	 
%----------------------------------------------------------------------------------------


%\newpage
%\loadgeometry{marginlayout}  
%\pagestyle{scrheadings} % Print headers again  
%\KOMAoptions{
%		headwidth=textwithmarginpar,
%		headsepline=true,
%		headsepline= 0.5pt,
%		footwidth=textwithmarginpar,
%	}   
%\onehalfspacing   
%
%
%\section{Équations rationnelles}
%
% 
% \begin{theorem}[Équation quotient nul] Si $\frac{A}{B}=0$ alors $A=0$ et $B\neq0$
% \end{theorem}
% \begin{proof} Comme $\frac{A}{B}=A\times \frac{1}{B}$, $B$ doit avoir un inverse. D'où $B\neq0$. 
% 	
% 	\noindent \parbox{0.85\linewidth}{	
% 		$\begin{WithArrows}
% 			\frac{A}{B} &=0 \\
% 			A\times \frac{1}{B}&=  0\Arrow{$\times B$}\\
% 			A\times \frac{1}{B}\times B & = 0\times B\\
% 			A & = 0
% 		\end{WithArrows}$ 	
% 	} 
% \end{proof}


%----------------------------------------------------------------------------------------
%	EXERCICES
%----------------------------------------------------------------------------------------

\newpage 
\loadgeometry{widelayout}  
\pagestyle{scrheadings} % Print headers again  
\KOMAoptions{
headwidth=textwithmarginpar,
headsepline=true,
headsepline= 0.5pt,
footwidth=textwithmarginpar,
} 
\onehalfspacing
\subsection{Exercices : résolution d'équations rationnelles}
\begin{example}Résoudre dans $\mathbb{R}$ les équations rationnelles suivantes

\noindent  \parbox{0.31\linewidth}{	
$\begin{WithArrows}
\frac{9x-5}{8x-3} &=   0\\
&  \rule{0pt}{1.3em} \\
&   \rule{0pt}{1.3em}  \\
&    \rule{0pt}{1.3em}  \\
&   \rule{0pt}{1.3em}   \\
&   \rule{0pt}{1.3em}   \\
&   \rule{0pt}{1.3em}   
\end{WithArrows}$ 
}  \hfill \parbox{0.31\linewidth}{	
$\begin{WithArrows}
\frac{-9x-5}{3x-4}	&= 5  \\
&  \rule{0pt}{1.3em} \\
&   \rule{0pt}{1.3em}  \\
&   \rule{0pt}{1.3em}  \\
&  \rule{0pt}{1.3em}   \\
&   \rule{0pt}{1.3em}    \\
&   \rule{0pt}{1.3em}   
\end{WithArrows}$ 
}    \hfill \parbox{0.31\linewidth}{	
$\begin{WithArrows} 
\frac{2}{6x-1}&= \frac{3}{5x-7}  \\
&  \rule{0pt}{1.3em} \\
&   \rule{0pt}{1.3em}  \\
&   \rule{0pt}{1.3em}  \\
&   \rule{0pt}{1.3em}   \\
&   \rule{0pt}{1.3em}    \\
&   \rule{0pt}{1.3em}   
\end{WithArrows}$ 
}  
\end{example}  
\vfill 
\begin{exercise}  \label{chapter08:exercice:equationrationnelles02}  
Préciser le domaine de résolution puis résoudre  dans $\mathbb{R}$ les équations suivantes : 
\begin{multicols}{4}
\begin{enumerate}[label={$(E_\arabic*)$\rule{0pt}{2em}},leftmargin=*]
\item $\frac{7x-6}{3x-5}=0$
\item $\frac{4x+6}{2x-5}=0$
\item $\frac{20x^2-8x}{2x-2}=0$
\item $\frac{6x^2-15x}{2x-5}=0$ 
\end{enumerate}
\end{multicols}  
\end{exercise} 
\begin{exercise}  \label{chapter08:exercice:equationrationnelles02b}  
Préciser le domaine de résolution puis résoudre  dans $\mathbb{R}$ les équations suivantes : 
\begin{multicols}{4}
\begin{enumerate}[label={$(E_\arabic*)$\rule{0pt}{2em}},leftmargin=*] 
\item $\frac{x+4}{-3x-3} =\frac{7}{2}$ % -29/23
\item $\frac{-x-6}{3-2x} =5$ 
\item $\frac{-6x-1}{-3x-3}=11$
\item $\frac{2x-6}{2x+1}=-11$
\end{enumerate}
\end{multicols}  
\end{exercise} 
 

\begin{exercise}   \label{chapter08:exercice:equationrationnelles03} 
Préciser le domaine de résolution, mettre au même dénominateur puis résoudre  dans $\mathbb{R}$ les équations suivantes  : 
\begin{multicols}{3}
\begin{enumerate}[label={$(E_\arabic*)$\rule{0pt}{2em}},leftmargin=*]
\item $1-\frac{2x}{2x+1}=\frac{3}{2x}$ 
\item $\frac{1}{x}-\frac{1}{x+1}=\frac{1}{x^2}$  
\item$x-\frac{1}{x}=0$ 
\item  $x-1=\frac{4}{x-1}$  
\item $1-\frac{3}{x-1}= \frac{x-4}{x-5} $ 
\item $\frac{3}{x+3}-\frac{15}{3x-5} =\frac{3}{2(x+3)}$ 
%\item $	\frac{7}{2x-1}  =\frac{3}{7x}$ 
\end{enumerate}
\end{multicols}   
\end{exercise}








%----------------------------------------------------------------------------------------
%	EXERCICES
%----------------------------------------------------------------------------------------

\newpage 

\begin{proof}[solutions de l'exercice \ref{chapter07:exercice:equationproduitA}]  
	 \scriptsize  
\begin{pycode}
from re import sub
import sympy as sp
from sympy.core.sympify import kernS

x, y, z, t , a, b = sp.symbols('x y z t a b', real=True)
k, m, n = sp.symbols('k m n', integer=True)
f, g, h = sp.symbols('f g h', cls=sp.Function)

# Create a list of functions to include in the table
funcs = [
' (5*x-10)*(8*x+5)' ,
' (8*x-1)*(7*x-4)' ,
' (8*x-1)+(7*x-4)' ,
' 5*x**2+2*x' ,
'  9*x**2-4' ,
' 36*x**2+100' ,
' x**2+10*x+25' ,
' -3*x**2-5*x' ,
' -x**2+2*x-1' ,
' x**2+4*x+4' ,
' (x-3)**2*(2*x-1)**4' ,
' x**2-6*x+9' ,
]
n=1 
for func in funcs: 
	solution = set(sp.solve(eval(func), x))
	string =  '$S_{'+str(n)+'}='+   sp.latex(solution ) + '$; '
	print(string) 
	n = n+1
\end{pycode}   
\end{proof}

\begin{proof}[solutions de l'exercice \ref{chapter07:exercice:equationproduitB}]  
	 \scriptsize  
\begin{pycode}
from re import sub
import sympy as sp
from sympy.core.sympify import kernS

x, y, z, t , a, b = sp.symbols('x y z t a b', real=True)
k, m, n = sp.symbols('k m n', integer=True)
f, g, h = sp.symbols('f g h', cls=sp.Function)

# Create a list of functions to include in the table
funcs = [
' (3*x+2)*(8*x+5)+(3*x+2)*(2*x+3)' ,
' (3*x-1)**2+(3*x-1)*(5*x-4)' ,
' (3*x-1)**2-(5*x-4)**2' ,
' (2*x-5)*(5*x-4)-(2*x-5)*(8*x-1)' ,
' (6*x-4)*(-3*x+2)-(10*x-2)*(6*x-4)' ,
' (2*x+3)**2-(5*x+4)**2' ,
]
n=1 
for func in funcs: 
	solution = set(sp.solve(eval(func), x))
	string =  '$S_{'+str(n)+'}='+   sp.latex(solution ) + '$; '
	print(string) 
	n = n+1
\end{pycode}  
\end{proof}
\begin{proof}[solution de l'exercice \ref{chapitre02:fractionalgebriques:simplification}] \hfill 
%\begin{multicols}{4} 
 
\begin{multicols}{2}  
\begin{pycode}
from re import sub
import sympy as sp

x, y, z, t = sp.symbols('x y z t')
k, m, n = sp.symbols('k m n', integer=True)
f, g, h = sp.symbols('f g h', cls=sp.Function)

# Create a list of functions to include in the table
funcsfac = [
#'(2*x-8)/(12*x-4)',
'(2*x**2+8*x)/(4*x)',
'(5*x-15)/(6*x-18)',
'(x**2-9)/(x+3)',
'(x-2)/(x**2-2*x)',
'(x**2-10*x+25)/(x-5)',
#'(x**2+4*x**2)/(x**2+x)',
#'(x**3+4*x**2)/(6*x**2+4*x)',
'(x+5)/(x**2-25)',
'(x**2+x)/(x**2-1)',
#'(x**2+x**2)/(3*x**3)',
'(9*x**2-12*x+4)/(9*x**2-4)'
]
n=0
L=['A','B','C','D','E','F', 'G', 'H', 'I', 'J', 'K', 'L', 'M']
for func in funcsfac :
	num, den =  sp.fraction(sp.together(func))
	roots = set(sp.solve(den, x))
	nums, dens = sp.cancel(sp.together(func)).as_numer_denom()
	myfactor =  eval('nums/dens')
	print( '$'+ str(L[n])  + "(x)=" + sp.latex(myfactor) + '$; pour $x\in\mathbb{R}\setminus'+ sp.latex(roots ) + '$; \n' )
	n     = n + 1
\end{pycode}
\end{multicols}
\end{proof}



\begin{proof}[solution de l'exercice \ref{chapitre02:fractionalgebriques:memedenominateur} ] \hfill 
\begin{multicols}{2}  
\begin{pycode} 
from re import sub
import sympy as sp

x, y, z, t = sp.symbols('x y z t')
k, m, n = sp.symbols('k m n', integer=True)
f, g, h = sp.symbols('f g h', cls=sp.Function)

# Create a list of functions to include in the table
funcsfac = [
'2/(x+3)+2/(x-3)',
'4/x+(x-1)/(3*x-5)',
'3/(5*x)+(11*x)/(x+1)',
'(3*x-2)/(5*x-3)-(2-3*x)/(7*x-2)',
'5/(2*x-1)-1',
'2*x+1-2/(2*x+1)',
'2/(x-2)-5/(x+3)',
'1/x-3/(x-5)',
'1/(2*x-7)-1/(2*x+9)',
#'(3*x+4)/(2*x-9)*(4*x**2-36*x+81)/(9*x**2-16)', 
#'(x**2-2*x+1)/(23*x-5)**2 *(46*x-10)/(1-x)',
]
n=0
L=['A','B','C','D','E','F', 'G', 'H', 'I', 'J', 'K', 'L', 'M']
for func in funcsfac :
	num, den =  sp.fraction(sp.together(func))
	roots = set(sp.solve(den, x))
	nums, dens = sp.cancel(sp.together(func)).as_numer_denom()
	myfactor =  eval('nums/dens')
	print( '$'+ str(L[n])  + "(x)=" + sp.latex(myfactor) + '$; pour $x\in\mathbb{R}\setminus'+ sp.latex(roots ) + '$; \n' )
	n     = n + 1
\end{pycode}
\end{multicols}
\end{proof}  
\begin{proof}[solution de l'exercice \ref{chapitre02:fractionalgebriques:produit} ] \hfill 
\begin{multicols}{2}  
\begin{pycode} 
from re import sub
import sympy as sp

x, y, z, t = sp.symbols('x y z t')
k, m, n = sp.symbols('k m n', integer=True)
f, g, h = sp.symbols('f g h', cls=sp.Function)

# Create a list of functions to include in the table
funcsfac = [ 
'(3*x+4)/(2*x-9)*(4*x**2-36*x+81)/(9*x**2-16)', 
'(x**2-2*x+1)/(23*x-5)**2 *(46*x-10)/(1-x)',
'1/((x+2)/(x-5))*(x+5)'
]
n=0
L=['A','B','C','D','E','F', 'G', 'H', 'I', 'J', 'K', 'L', 'M']
for func in funcsfac :
	num, den =  sp.fraction(sp.together(func))
	roots = set(sp.solve(den, x))
	nums, dens = sp.cancel(sp.together(func)).as_numer_denom()
	myfactor =  eval('nums/dens')
	print( '$'+ str(L[n])  + "(x)=" + sp.latex(myfactor) + '$; pour $x\in\mathbb{R}\setminus'+ sp.latex(roots ) + '$; \n' )
	n     = n + 1
\end{pycode}
\end{multicols}
\end{proof}  



\begin{multicols}{3}
\begin{proof}[solution de l'exercice  \ref{chapter08:exercice:equationrationnelles02} ] \hfill \footnotesize

\begin{pycode}
from re import sub
import sympy as sp

x, y, z, t = sp.symbols('x y z t')
k, m, n = sp.symbols('k m n', integer=True)
f, g, h = sp.symbols('f g h', cls=sp.Function)

# Create a list of functions to include in the table
funcsfac = [
'(7*x-6)/(3*x-5)',
'(4*x+6)/(2*x-5)',
'(20*x**2-8*x)/(2*x-2)',
'(6*x**2-15*x)/(2*x-5)',
]
n=0
L=['S_1','S_2','S_3','S_4','S_5','S_6', 'S_7', 'S_8', 'S_9', 'S_{10}', 'S_{11}', 'S_{12}']
for func in funcsfac :
	num, den = sp.fraction(func)
	roots = set(sp.solve(den, x))
	string =  "\\noindent$(E_"+ str(n+1) + ")$  D.R.  $=\mathbb{R}\setminus" + sp.latex(roots ) + '$;   '
	print(string)
	solution = set(sp.solve(func, x))
	string =  '$'+ str(L[n])  + '=' + sp.latex(solution ) + '$; \n '
	print(string)
\end{pycode}
\end{proof}
\begin{proof}[solution de l'exercice \ref{chapter08:exercice:equationrationnelles02b} ] \hfill \footnotesize

\begin{pycode}
from re import sub
import sympy as sp

x, y, z, t = sp.symbols('x y z t')
k, m, n = sp.symbols('k m n', integer=True)
f, g, h = sp.symbols('f g h', cls=sp.Function)

# Create a list of functions to include in the table
funcsfac = [
'(x+4)/(-3*x-3)-7/2',
'(-x-6)/(3-2*x)-5',
'(-6*x-1)/(-3*x-3)-11',
'(2*x-6)/(2*x+1)+11',
]
n=0
L=['S_1','S_2','S_3','S_4','S_5','S_6', 'S_7', 'S_8', 'S_9', 'S_{10}', 'S_{11}', 'S_{12}']
for func in funcsfac :
	num, den = sp.together(func).as_numer_denom()
	roots = set(sp.solve(den, x))
	string =  "\\noindent$(E_"+ str(n+1) + ")$  D.R.  $=\mathbb{R}\setminus" + sp.latex(roots ) + '$;   ' 
	print(string)
	solution = set(sp.solve(func, x))
	string =  '$'+ str(L[n])  + '=' + sp.latex(solution ) + '$; \n '
	print(string)
	n     = n + 1
\end{pycode}
\end{proof} 
 
\begin{proof}[solution de l'exercices \ref{chapter08:exercice:equationrationnelles03} ] \hfill \footnotesize

\begin{pycode}
from re import sub
import sympy as sp

x, y, z, t = sp.symbols('x y z t')
k, m, n = sp.symbols('k m n', integer=True)
f, g, h = sp.symbols('f g h', cls=sp.Function)

# Create a list of functions to include in the table
funcsfac = [ 
'1-(2*x)/(2*x+1)-3/(2*x)',
'1/x-1/(x+1)-1/(x**2)',
'x-1/x',
'x-1-4/(x-1)',
'1-3/(x-1)-(x-4)/(x-5)',
'3/(x+3)-15/(3*x-5)-3/(2*x+6)',
#'-2+(5*x-4)/(2*x-5)-(x+6)/(x-5)',
#'(25*x)/(x+2)-25/(3*x+1)+2/((3*x+1)*(x+2))' 
]
n=0
L=['S_1','S_2','S_3','S_4','S_5','S_6', 'S_7', 'S_8', 'S_9', 'S_{10}', 'S_{11}', 'S_{12}']
for func in funcsfac :
	num, den = sp.together(func).as_numer_denom()
	roots = set(sp.solve(den, x))
	string =  "\\noindent$(E_"+ str(n+1) + ")$  D.R.  $=\mathbb{R}\setminus" + sp.latex(roots ) + '$;   ' 
	print(string)
	solution = set(sp.solve(func, x))
	string =  '$'+ str(L[n])  + '=' + sp.latex(solution ) + '$; \n '
	print(string)
	n     = n + 1
\end{pycode}
\end{proof} 
\end{multicols} 
 

%----------------------------------------------------------------------------------------
%	
%----------------------------------------------------------------------------------------
\newpage  
\section{Club de Maths : équations simples et moins simples}


\begin{problem}
\label{chapter07:problem:01}\hfill 

\noindent\parbox{0.35\linewidth}{\centering 
\begin{tikzpicture}[scale=0.6]
\tkzSetUpPoint[size=2, fill=black]
\tkzfctset{tan style/.style={->,>=latex, color=red, line width=1.2pt}}   
\tkzSetUpLine[line width= 1.2pt, add=.5 and .5]  
%			\tkzInit[xmin=-1.0,ymin=-1.0,xmax=6.5,ymax=2.8]
%			\tkzClip
\tkzDefPoint(0,0){B} 
\tkzDefPoint(6,0){C}
\tkzDefPoint(6,1.25){D}
\tkzDefPoint(2,1.25){E}
\tkzDefPoint(2,4.2){F}
\tkzDefPoint(0,4.2){A}
\tkzDrawPolygon[fill=gray!15, line width=1.2pt](A,B,C,D,E,F)  

\tkzDrawSegment[dim={\scriptsize$10$,-15pt,}](D,E))
\tkzDrawSegment[dim={\scriptsize$5$,+10pt,}](A,F))
\tkzDrawSegment[dim={\scriptsize$2$,-10pt,},dashed](C,D))
\tkzDrawSegment[line width=1.2pt](A,C)


\tkzLabelPoints[above right,font=\scriptsize](D,F)
\tkzLabelPoints[above left,font=\scriptsize](A)
\tkzLabelPoints[below right,font=\scriptsize](C)
\tkzLabelPoints[below left,font=\scriptsize](B,E)


\tkzMarkRightAngle[  size=0.25](B,A,F)
\tkzMarkRightAngle[  size=0.25](C,B,A)
\tkzMarkRightAngle[  size=0.25](D,C,B)
\tkzMarkRightAngle[  size=0.25](C,D,E)
\tkzMarkRightAngle[  size=0.25](F,E,D)
\tkzMarkRightAngle[  size=0.25](A,F,E)
\end{tikzpicture}  
}\hfill \parbox{0.6\linewidth}{ 
Quelle doit être la longueur $EF$ pour que l'aire de la partie colorée soit égale à l'aire du triangle $ABC$ ?
}    

\end{problem}

\begin{problem}\label{chapter07:problem:02} 
Un père et son fils ont $35$ et $3$ ans, respectivement. Dans combien d'années le père aura-t-il le double de l'âge de son fils ?
% 35+x=2(3+x) , x=29$.
\end{problem}
\begin{problem}\label{chapter07:problem:03} 
Un joueur obtient huit fois le score $7$ puis $x$ fois le score $10$. À la fin du jeu, son score moyen est de $9$. Combien vaut alors $x$ ?

% $\frac{8\times 7 +10x}{x+8}=9$, $x=16$.
\end{problem}
\begin{problem}\label{chapter07:problem:04}  
Résoudre dans $\mathbb{R}$ en développant et en simplifiant les équations suivantes d'inconnue $x$ :\vspace{-0.5em}
\begin{multicols}{2}
\begin{enumerate}[label={$(E_\arabic*)$},leftmargin=*]
\item $(x-1)^2+(x+3)^2=2(x-2)(x+1)+38$
\item $5(x^2-2x-1)+2(3x-2)=5(x+1)^2$ 
\end{enumerate}
\end{multicols} 
\end{problem}
\begin{problem}\label{chapter07:problem:05} 
Résoudre dans $\mathbb{R}$ en factorisant   les équations suivantes d'inconnue $x$ :\vspace{-0.5em}
\begin{multicols}{2}
\begin{enumerate}[label={$(E_\arabic*)$},leftmargin=*]
\item $-\frac{3x^2}{5}+x=0$
\item $-\frac{5x^2}{7}-\frac{3x}{4}=0$ 
\end{enumerate}
\end{multicols} 
\end{problem}
\begin{problem}\label{chapter07:problem:06} 
Résoudre dans $\mathbb{R}$ en factorisant   les équations suivantes d'inconnue $x$ :\vspace{-0.5em}
\begin{multicols}{2}
\begin{enumerate}[label={$(E_\arabic*)$},leftmargin=*]
\item $(x+5)(4x-1)+x^2-25=0$
\item $(x+4)(5x+9)-x^2+16=0$ 
\end{enumerate}
\end{multicols} 
\end{problem}
\begin{problem}\label{chapter07:problem:07} 
Résoudre dans $\mathbb{R}$ en factorisant   les équations suivantes d'inconnue $x$ :\vspace{-0.5em}
\begin{multicols}{2}
\begin{enumerate}[label={$(E_\arabic*)$},leftmargin=*]
\item $7x^3-175x=0$
\item $(x+5)(3x+2)^2=x^2(x+5)$ 
\end{enumerate}
\end{multicols} 
\end{problem}
\begin{problem}\label{chapter07:problem:08} 
Résoudre dans $\mathbb{R}$   les équations suivantes d'inconnue $x$ :\vspace{-0.5em}
\begin{multicols}{3}
\begin{enumerate}[label={$(E_\arabic*)$},leftmargin=*]
\item $\frac{4x+7}{x-1}=\frac{12x+5}{3x+4}$
\item $\frac{7}{x-5}=\frac{4}{x+1} +\frac{3}{x-2}$ 
\item $\frac{9}{x}=\frac{8}{x+1}+\frac{1}{x-1}$ 
\end{enumerate}
\end{multicols} 
\end{problem}
\begin{problem}\label{chapter07:problem:09}  
On note $a$ une longueur, pour l’instant inconnue. Un rectangle de largeur $a$ et de longueur
$a+1$ a la même aire qu’un autre rectangle qui lui est de largeur $a−1$ et de longueur $a+3$. Est-ce possible et si oui, combien de valeurs possibles de $a$ existe-t-il ?
\end{problem}
\begin{problem}\label{chapter07:problem:10} 
Une boîte de café coûte \EUR{100} à l'achat et se vend \EUR{140} : elle permet une marge de 40\%.\\
Pour une boîte de cacao, la marge est de 20\%. Si le nombre de boîtes de café vendues est le double du nombre de boîte de cacao vendues, et si la marge totale est de 36\% (par rapport au prix d'achat), à quel prix s'est vendue chaque boîte de cacao ? 
\end{problem} 
\begin{problem}\label{chapter07:problem:11} 
Si $\frac{a}{b}=\frac{1}{2}$ et $\frac{b}{c}=\frac{1}{4}$, donner la valeur de $\frac{a-b}{b-c}$.

% $\frac{1}{6} $ 
\end{problem} 
\begin{problem}\label{chapter07:problem:12} 
Trouver $x$ sachant que $\frac{5}{9}=\dfrac{1}{1+\dfrac{1}{1+\dfrac{1}{x}}}$.
\end{problem}
%----------------------------------------------------------------------------------------
%	
%----------------------------------------------------------------------------------------
%\newpage  
\begin{remark} Ces exercices sont pour une grande partie tirés d'un recueil du club de Maths de Nancy.
	\end{remark}
\begin{proof}[solution du problème \ref{chapter07:problem:01}]  

$x=EF$. $7.5(x+2)=5x+30$. $x=6$.  

\end{proof} 
\begin{proof}[solution du problème \ref{chapter07:problem:02}]  

$35+x=2(3+x)$ , $x=29$.  

\end{proof} 
\begin{proof}[solution du problème \ref{chapter07:problem:03}]  

$\frac{8\times 7 +10x}{x+8}=9$, $x=16$.
\end{proof} 
\begin{proof}[solution du problème \ref{chapter07:problem:04}]  

Pour vérification, chacune des deux équations admet une solution, et leur somme faut $3$.

\end{proof} 
\begin{proof}[solution du problème \ref{chapter07:problem:05}]  

Pour vérification : chacune des deux équations admet deux solutions et la somme des quatre solutions vaut $\frac{37}{60}$.

\end{proof} 
\begin{proof}[solution du problème \ref{chapter07:problem:06}]  

Pour vérification : chacune des deux équations admet deux solutions et la somme des quatre solutions vaut $\frac{-171}{20}$.

\end{proof} 
\begin{proof}[solution du problème \ref{chapter07:problem:07}]  

Pour vérification : chacune des deux équations admet trois solutions et la somme des six solutions vaut $\frac{-13}{2}$.

\end{proof}  
\begin{proof}[solution du problème \ref{chapter07:problem:08}]  

Pour vérification : chacune des trois équations admet une solution et la somme des trois solutions vaut $\frac{529}{308}$.

\end{proof} 
\begin{proof}[solution du problème \ref{chapter07:problem:09}]  

Les termes quadratiques (en $a^2$) se simplifient et $a=3$.

\end{proof} 
\begin{proof}[solution du problème \ref{chapter07:problem:10}]  

$x=$ prix d'achat de la boîte de cacao, $n=$ nombre de boîtes de cacao vendues.  $0.20\times nx+2n\times 40 =0.36\times \underbrace{(nx+2n\times 100)}_{\text{prix d'achat total} }$.\\
$x$ est donc solution de $0.20x+80=0.36(x+200)$. $x=50$. 
Le prix de vente est de \EUR{60} .

\end{proof} 
 
%----------------------------------------------------------------------------------------
%	 
%----------------------------------------------------------------------------------------


	
%----------------------------------------------------------------------------------------
%	 
%----------------------------------------------------------------------------------------

\cleardoublepage 

%----------------------------------------------------------------------------------------
%	 
%----------------------------------------------------------------------------------------
\end{document}